\section{Method}
\label{sec:method}

\subsection{Contextual W/R/V attention}

Let $\mH\in\mathbb{R}^{B\times T\times d}$ be normalized hidden states. Each layer projects contextual reads, writes, and values,
\begin{align}
    \mR &= \mH\mW_r, &
    \mW_{\mathrm{ctx}} &= \mH\mW_w, &
    \mV &= \mH\mW_v.
\end{align}
The evaluated model uses eight read heads and two shared write/value heads, with head dimension 48. A deterministic token code and a learned projection of the tied rank-16 lexical table form a lexical write address $\mL(x_t)$. For write head $h$,
\begin{equation}
    \vw_{t,h}=\vw^{\mathrm{ctx}}_{t,h}
      +\tanh(g_h)\,\mL_h(x_t),
    \qquad g_h=0\ \text{at initialization}.
\end{equation}
Thus the lexical branch is exactly neutral initially; the contextual write stream remains fully active. Reads and writes receive independent RMS normalization per head and rotary position encoding. Write/value heads are repeated across their four read-head groups, and causal attention is
\begin{equation}
    \vy_{t,h}=\sum_{j\le t}
    \operatorname{softmax}_j\!\left(
      \frac{\operatorname{RoPE}(\vr_{t,h})^\top
      \operatorname{RoPE}(\vw_{j,\kappa(h)})}{\sqrt{48}}
    \right)\vv_{j,\kappa(h)},
\end{equation}
where $\kappa(h)$ maps each read head to one of the two shared write/value heads. PyTorch scaled dot-product attention is configured to permit FlashAttention dispatch on H100. W/R/V is therefore an attention mechanism; the terminology describes data flow and the lexical write residual, not the absence of softmax attention. Incremental execution retains past W and V tensors and grows with prefix length.

\subsection{Tied lexical residual object}

The shared feed-forward transformation is a dense SwiGLU,
\begin{equation}
    \mF(\vh)=\mW_d\left[\operatorname{SiLU}(\mW_g\vh)\odot \mW_u\vh\right].
\end{equation}
For token identity $x_t$, a tied table $\mE\in\mathbb{R}^{|\mathcal V|\times r}$ provides a low-rank modulation:
\begin{align}
    \vz_t &= \operatorname{SiLU}(\mU\vh_t)\odot(\mathbf{1}+\mE[x_t]),\\
    \mO(\vh_t,x_t) &= \mD\vz_t.
\end{align}
The same table $\mE$ is shared across layers; projection matrices remain layer-specific. Setting $\mE$ to zero makes the token-dependent modulation neutral initially, but the object branch itself remains active through a nonzero output gate $\alpha$:
\begin{equation}
    \mF_{\mathrm{object}}(\vh_t,x_t)=\mF(\vh_t)+\alpha\mO(\vh_t,x_t).
\end{equation}

\subsection{Deterministic micro-expert residual}

Each layer assigns token $x_t$ to one of $N$ micro-experts using
\begin{equation}
    e_\ell(x_t)=(x_t+\ell)\bmod N.
\end{equation}
All expert hidden activations are computed as one regular tensor operation and the selected output is gathered without sorting or scatter. With expert width $w$,
\begin{equation}
    \mR_\ell(\vh_t,x_t)=\beta\,\mW^{(e_\ell(x_t))}_{d}
    \left[\operatorname{SiLU}(\mW^{(e_\ell(x_t))}_{g}\vh_t)
    \odot \mW^{(e_\ell(x_t))}_{u}\vh_t\right].
\end{equation}
The complete feed-forward output is $\mF+\alpha\mO+\beta\mR$. This is a lexical residual over a shared dense path, not a claim that routing replaces shared computation.

\subsection{Fixed-state WVR associative ablation}

The earlier fixed-state ablation uses one recurrent associative matrix rather than softmax attention. We denote that mechanism by \emph{WVR}: write address, contextual value, and read address. For token $x_t$ and hidden state $\vh_t$, a compressed lexical forge reuses the tied lexical object table,
\begin{equation}
    \vf_t=\mW_{f,2}\operatorname{SiLU}(\mW_{f,1}\mE[x_t]),
\end{equation}
where the implemented bottleneck is $16\rightarrow4\rightarrow128$. The common address combines this deterministic lexical code with a projected hidden-state signature,
\begin{align}
    \vc_t &= \operatorname{norm}(\mW_c\vh_t),\\
    \va_t &= \operatorname{norm}\!\left(\operatorname{forge}(x_t)+
        \tanh(\eta)\vc_t\right).
\end{align}
The lexical code is an address tied to token identity; it is not evidence of semantic representation. The contextual projection is likewise described operationally rather than semantically.

Writing uses the preceding address and the current contextual payload,
\begin{align}
    \mathbf{w}_t &= \va_{t-1}, & \mathbf{v}_t &= \mW_v\vh_t, & \mathbf{r}_t &= \va_t,\\
    \widehat{\mathbf{v}}_t &= \mathbf{w}_t^\top\mM_{t-1},\\
    \mM_t &= \lambda\mM_{t-1}+\rho\mathbf{w}_t
        (\mathbf{v}_t-\widehat{\mathbf{v}}_t)^\top,\\
    \vy_t &= \mathbf{r}_t^\top\mM_{t-1},\\
    \vc^{\mathrm{WVR}}_t &= \gamma\mW_o\vy_t.
\end{align}
The context residual $\vc^{\mathrm{WVR}}_t$ is added to the causal-convolution output. The implementation shifts the write stream causally: the current read uses $\mathbf{r}_t$, while the preceding write address is paired with the current payload. The chunkwise triangular solve is an exact parallelization of this recurrence. Thus WVR stores a contextual hidden value rather than a lexical code. The state is one rank-$128$ matrix plus fixed-size address and collision-load vectors, independent of prefix length.

Collision normalization maintains a decayed diagonal load estimate $\vd_t$ and scales write and read coordinates by $(\vd_t+\epsilon)^{-1/2}$ before normalization. It does not partition the rank-$128$ matrix into separate memories.

QKV attention compares a contextual query with every retained key and performs normalized competition over individual token occurrences. WVR instead reads the pre-update compressed recurrent association through $\mathbf{r}_t^\top\mM_{t-1}$. Our training comparison concerns modeling quality and throughput without assigning credit for KV-cache size. The representational tradeoff is explicit: WVR has fixed state, but repeated or colliding occurrences may overwrite one another.

The experimental occurrence-address extension replaces the dense contextual projection by a compact $d\rightarrow8\rightarrow128$ signature and adds a gated deterministic position signature,
\begin{equation}
    \va_t=\operatorname{norm}\!\left(
    \operatorname{forge}(x_t)+\operatorname{smallctx}(\vh_t)
    +\tanh(\gamma)\operatorname{pos}(t)\right),
\end{equation}
with $\gamma=0$ initially and one fixed-size position counter for incremental decoding. Its completed full-budget run is reported as a negative ablation.

\subsection{Architectural path ablations}

We retain three completed intermediate conditions: one W/R/V layer inserted into a collision-normalized WVR model; eight W/R/V layers with strongly lexical normalized addresses; and eight contextual W/R/V layers with two remaining WVR layers. These are path ablations, not isolated factorial controls. Their ordering motivates hypotheses about insufficient selective attention, overly dominant lexical addressing, and residual fixed-state bottlenecks, but does not identify those causes independently.
